\chapter{Generating Functions as a Computing Tool}

\begin{center}\itshape
A generating function is a clothesline on which we hang up a sequence of numbers for display.\\
\upshape --- Herbert Wilf, \emph{generatingfunctionology}
\end{center}

Tome~3, Chapter~3, used generating functions of probability distributions: $\mathbb E[s^X]$ packs a whole law
into one function, and sums of independent variables become products. This chapter uses the same device as a pure
computing tool. A sequence $(a_n)$ becomes the series $\sum a_nx^n$, and operations on sequences that are awkward
(convolutions, recurrences, sums over all partitions of a set) become operations on functions that are easy
(products, quotients, exponentials). One can then read off a closed form, an exact recurrence, or, as Chapter~11
will show, the asymptotic size of $a_n$.

We first set up \emph{formal} power series, which make these manipulations rigorous without any question of
convergence. Then come three classical uses: linear recurrences, which are exactly the rational generating
functions; \emph{exponential} generating functions, which count labelled structures and obey the exponential
formula; and partitions of integers, with Euler's pentagonal number theorem, whose proof by Franklin is one of
the most charming bijections in mathematics.

\begin{think}
Expand $\frac1{1-x-x^2}=1+x+2x^2+3x^3+5x^4+\cdots$. Where do the Fibonacci numbers come from, and what does the
factorisation of $1-x-x^2$ tell you about them?
\end{think}

\section{Formal power series}

Let $K$ be $\Q$, $\R$ or $\C$. A \emph{formal power series} over $K$ is a sequence $(a_n)_{n\ge0}$ of elements of $K$,
written $A(x)=\sum_{n\ge0}a_nx^n$; nothing is assumed about convergence, and $x$ is not a number. We write
$[x^n]A=a_n$. Sum and product are defined by
\[
\begin{gathered}\sum a_nx^n+\sum b_nx^n=\sum(a_n+b_n)x^n,\\\Bigl(\sum a_nx^n\Bigr)\Bigl(\sum b_nx^n\Bigr)=\sum_n\Bigl(\sum_{k=0}^na_kb_{n-k}\Bigr)x^n.\end{gathered}
\]
With these operations the formal power series form a commutative ring $K[[x]]$ containing the polynomials $K[x]$:
associativity and distributivity are identities between finite sums of coefficients, checked as for polynomials.
The \emph{order} of a non-zero series is the least $n$ with $a_n\ne0$. Since $[x^n](AB)$ involves only $a_0,\dots,a_n$ and
$b_0,\dots,b_n$, the order of $AB$ is the sum of the orders (the lowest coefficients multiply to a non-zero number), so
$K[[x]]$ has no zero divisors.

\begin{proposition}\label{prop:inverse}
A series $A$ is invertible in $K[[x]]$ if and only if $a_0\ne0$. For instance $\frac1{1-cx}=\sum_nc^nx^n$, and more
generally, for $m\ge1$,
\[
\frac1{(1-cx)^m}=\sum_{n\ge0}\binom{n+m-1}{m-1}c^nx^n.
\]
\end{proposition}
\begin{proof}
If $AB=1$ then $a_0b_0=1$. Conversely, if $a_0\ne0$, the equations $\sum_{k=0}^na_kb_{n-k}=[n=0]$ are solved by
$b_0=\frac1{a_0}$ and $b_n=-\frac1{a_0}\sum_{k=1}^na_kb_{n-k}$, recursively. Then $(1-cx)\sum c^nx^n=1$ because the sum
telescopes. The coefficient of $x^n$ in $\bigl(\sum_jc^jx^j\bigr)^m$ is $c^n$ times the number of $m$-tuples of
non-negative integers with sum $n$, which is $\binom{n+m-1}{m-1}$ by the stars-and-bars count (Tome~0, Chapter~8).
\end{proof}

Infinite sums and products are allowed when they make sense coefficientwise: if $A_1,A_2,\dots$ are series whose
orders tend to infinity, each coefficient of $\sum_kA_k$ involves only finitely many $A_k$, so the sum is defined, and
likewise $\prod_k(1+A_k)$, where each coefficient involves the finitely many factors with order at most $n$. In
particular, if $G$ has $g_0=0$, the \emph{composition} $F(G)=\sum_kf_kG^k$ is defined, since $G^k$ has order at least $k$.

The \emph{derivative} is $DA=\sum_{n\ge1}na_nx^{n-1}$. It is linear and satisfies the product rule $D(AB)=(DA)B+A(DB)$,
since both sides have coefficient $\sum_k(k+1)a_{k+1}b_{n-k}+\sum_ka_k(n-k+1)b_{n-k+1}=(n+1)\sum_ka_kb_{n+1-k}$ at
$x^n$. In characteristic $0$, $DA=0$ forces $A$ to be constant, and the chain rule holds:
\[
D\bigl(F(G)\bigr)=(DF)(G)\cdot DG\qquad(g_0=0).
\]
Indeed, for $F=x^k$ it is the product rule iterated, $D(G^k)=kG^{k-1}DG$; by linearity it holds for polynomials; and the
coefficient of $x^n$ on either side only involves the coefficients $f_0,\dots,f_{n+1}$ of $F$, so it holds for $F$
if it holds for the polynomial $f_0+\dots+f_{n+1}x^{n+1}$.

We define $\exp(x)=e^x=\sum_n\frac{x^n}{n!}$ and $\log(1+x)=\sum_{n\ge1}\frac{(-1)^{n-1}x^n}n$, so that $D\exp=\exp$ and
$D\log(1+x)=\frac1{1+x}$.

\begin{proposition}\label{prop:explog}
Let $F$, $G$ be series with $f_0=g_0=0$. Then $\exp(F+G)=\exp(F)\exp(G)$, $\log(1+(e^F-1))=F$, and
$\exp\bigl(\log(1+F)\bigr)=1+F$.
\end{proposition}
\begin{proof}
By the chain rule $D\exp(F)=\exp(F)DF$. The series $H=\exp(F+G)\exp(-F)\exp(-G)$ has
$DH=H\cdot D(F+G-F-G)=0$ by the product rule, so $H$ is constant, equal to its constant term $1$; and
$\exp(-F)\exp(F)=1$ by the same argument with $G=-F$. This gives the first identity. For the second, put
$U=e^F-1$: then $D\log(1+U)=\frac{DU}{1+U}=\frac{e^FDF}{e^F}=DF$, and both sides vanish at $0$. For the third, let
$L=\log(1+F)$, so that $DL=\frac{DF}{1+F}$; then $D\bigl(e^L(1+F)^{-1}\bigr)=e^L(1+F)^{-1}\bigl(DL-\frac{DF}{1+F}\bigr)=0$,
and the constant term is $1$.
\end{proof}

\section{Two tools from analysis}

Generating functions are most useful when they also converge, and then two facts about infinite sums are needed again
and again in this tome. Tome~0 used them only in special cases, so we prove them here.

For a family $(a_i)_{i\in I}$ of non-negative numbers indexed by a countable set $I$, let $\sum_{i\in I}a_i\in[0,+\infty]$ be the
supremum of the finite sums $\sum_{i\in F}a_i$ over finite $F\subseteq I$. A complex family is \emph{summable} if
$\sum_{i\in I}|a_i|<\infty$.

\begin{lemma}[summation by packets]\label{lem:packets}
\begin{enumerate}[label=(\alph*)]
\item If $a_i\ge0$, then for every enumeration $i_1,i_2,\dots$ of $I$, $\sum_{k\ge1}a_{i_k}=\sum_{i\in I}a_i$.
\item If $a_i\ge0$ and $I$ is the disjoint union of countably many sets $I_j$, then $\sum_{i\in I}a_i=\sum_j\sum_{i\in I_j}a_i$.
\item If $(a_i)$ is a summable complex family, the series $\sum_ka_{i_k}$ converges absolutely to the same sum $S$ for every
enumeration, and $S=\sum_j\sum_{i\in I_j}a_i$ for every partition into packets $I_j$, each inner sum being itself absolutely convergent.
In particular, for a double sequence with $\sum_{m,n}|a_{mn}|<\infty$, $\sum_m\sum_na_{mn}=\sum_n\sum_ma_{mn}$.
\end{enumerate}
\end{lemma}
\begin{proof}
(a) Each partial sum $\sum_{k\le K}a_{i_k}$ is a finite sum, hence at most $\sum_Ia_i$; and each finite $F\subseteq I$ is contained in
$\{i_1,\dots,i_K\}$ for some $K$, so $\sum_Fa_i\le\sum_{k\le K}a_{i_k}$. The two suprema are equal.

(b) A finite $F\subseteq I$ meets finitely many $I_j$, and $\sum_Fa_i=\sum_j\sum_{F\cap I_j}a_i\le\sum_j\sum_{I_j}a_i$; so
$\sum_Ia_i\le\sum_j\sum_{I_j}a_i$. Conversely, for finitely many indices $j\le J$ and finite sets $F_j\subseteq I_j$,
$\sum_{j\le J}\sum_{F_j}a_i=\sum_{\bigcup F_j}a_i\le\sum_Ia_i$; taking suprema over the $F_j$, then over $J$, gives the reverse
inequality.

(c) Write $a_i=b_i-c_i+i(d_i-e_i)$ with $b_i=\max(\operatorname{Re}a_i,0)$, $c_i=\max(-\operatorname{Re}a_i,0)$ and similarly $d_i$, $e_i$
for the imaginary part. These four families are non-negative with sums at most $\sum|a_i|<\infty$, so (a) and (b) apply to each,
with finite values, and the statements for $(a_i)$ follow by linearity of convergent series. Absolute convergence of each
enumeration and of each inner sum holds by (a) and (b) applied to $(|a_i|)$.
\end{proof}

A sequence of functions $f_n$ on a set $E$ \emph{converges uniformly} to $f$ if $\sup_E|f_n-f|\to0$.

\begin{lemma}[uniform convergence]\label{lem:uniform}
Let $E$ be an interval.
\begin{enumerate}[label=(\alph*)]
\item A uniform limit of continuous functions on $E$ is continuous.
\item \emph{Weierstrass test}: if $|u_n(x)|\le M_n$ on $E$ with $\sum M_n<\infty$, then $\sum u_n$ converges uniformly on $E$.
\item If $f_n\to f$ uniformly on $[a,b]$ and the $f_n$ are continuous, then $\int_a^bf_n\to\int_a^bf$. So a uniformly convergent
series of continuous functions can be integrated term by term.
\item If the $f_n$ are of class $C^1$ on $E$, $f_n(x_0)$ converges for some $x_0\in E$, and $f_n'\to g$ uniformly on $E$, then $f_n$
converges, uniformly on every bounded subinterval, to a function $f$ of class $C^1$ with $f'=g$. So a series of $C^1$ functions can
be differentiated term by term when the differentiated series converges uniformly.
\end{enumerate}
\end{lemma}
\begin{proof}
(a) Let $x\in E$ and $\varepsilon>0$. Choose $n$ with $\sup|f_n-f|\le\varepsilon$, then $\delta>0$ with $|f_n(y)-f_n(x)|\le\varepsilon$ for
$|y-x|\le\delta$. Then $|f(y)-f(x)|\le|f(y)-f_n(y)|+|f_n(y)-f_n(x)|+|f_n(x)-f(x)|\le3\varepsilon$.
(b) The series converges absolutely at each $x$, and the tail satisfies $\bigl|\sum_{k>n}u_k(x)\bigr|\le\sum_{k>n}M_k\to0$ uniformly in $x$.
(c) $f$ is continuous by (a), hence integrable, and $\bigl|\int_a^b(f_n-f)\bigr|\le(b-a)\sup|f_n-f|\to0$.
(d) By (a), $g$ is continuous. By the fundamental theorem of calculus, $f_n(x)=f_n(x_0)+\int_{x_0}^xf_n'$. By (c) the right side
converges to $f(x)=\ell+\int_{x_0}^xg$, where $\ell=\lim f_n(x_0)$, uniformly for $x$ in a bounded subinterval $J$, since
$\bigl|\int_{x_0}^x(f_n'-g)\bigr|\le|x-x_0|\sup_E|f_n'-g|$. And $f'=g$ by the fundamental theorem again.
\end{proof}

\begin{remark}
When the series converge, formal computations agree with the functions. If $\sum a_nz^n$ and $\sum b_nz^n$ converge
absolutely for $|z|<r$, their Cauchy product converges to the product of the sums (Tome~0, Chapter~5). And a series
with positive radius of convergence is determined by its sum near $0$: if $\sum c_nz^n=0$ for all small $z$ and $c_m$ is
the first non-zero coefficient, dividing by $z^m$ and letting $z\to0$ gives $c_m=0$. So an identity between functions
that are sums of power series near $0$ gives an identity of formal series, and conversely. We use this freely below.
\end{remark}

\section{Rational series and linear recurrences}

\begin{theorem}\label{thm:rational}
Let $a_1,\dots,a_d\in\C$ with $a_d\ne0$, and $Q(x)=1-a_1x-\dots-a_dx^d$. For a sequence $(u_n)$ with generating
function $U=\sum u_nx^n$, the following are equivalent:
\begin{enumerate}[label=(\roman*)]
\item $u_{n+d}=a_1u_{n+d-1}+\dots+a_du_n$ for every $n\ge0$;
\item $U=\frac PQ$ for a polynomial $P$ of degree less than $d$.
\end{enumerate}
In that case, if $Q(x)=\prod_{i=1}^s(1-\alpha_ix)^{m_i}$ with distinct $\alpha_i$, there are polynomials $p_i$ of degree less
than $m_i$ such that $u_n=\sum_{i=1}^sp_i(n)\alpha_i^n$ for all $n\ge0$.
\end{theorem}
\begin{proof}
The coefficient of $x^{n+d}$ in $QU$ is $u_{n+d}-a_1u_{n+d-1}-\dots-a_du_n$. So (i) says exactly that $P=QU$ has no
terms of degree $\ge d$, which is (ii) since $Q$ is invertible.

For the explicit form, note first that $Q(0)=1$ and that $Q$ factors over $\C$ as $\prod(1-\alpha_ix)^{m_i}$: the $\alpha_i$
are the inverses of the roots of $Q$ (equivalently, the roots of $y^d-a_1y^{d-1}-\dots-a_d$, all non-zero since $a_d\ne0$).
We show by induction on the number $s$ of distinct factors that $\frac PQ$, with $\deg P<\deg Q$, is a sum
$\sum_i\frac{P_i}{(1-\alpha_ix)^{m_i}}$ with $\deg P_i<m_i$. For $s=1$ there is nothing to do. For $s\ge2$, write $Q=Q_1Q_2$
with $Q_1=(1-\alpha_1x)^{m_1}$; $Q_1$ and $Q_2$ have no common root, so they are coprime, and since $\C[x]$ is a principal
ideal domain (Tome~2, Chapter~6) there are polynomials $A$, $B$ with $AQ_1+BQ_2=1$. Then
$\frac PQ=\frac{PA}{Q_2}+\frac{PB}{Q_1}$. Divide $PB$ by $Q_1$: $PB=CQ_1+R_1$ with $\deg R_1<m_1$, so
$\frac PQ=\frac{R_1}{Q_1}+\frac{PA+CQ_2}{Q_2}$. Write $R_2=PA+CQ_2$. Multiplying by $Q$ gives $P=R_1Q_2+R_2Q_1$, and since
$\deg P<\deg Q$ and $\deg(R_1Q_2)<\deg Q$, also $\deg(R_2Q_1)<\deg Q$, so $\deg R_2<\deg Q_2$. The induction hypothesis
applies to $\frac{R_2}{Q_2}$.

Finally, a polynomial $P_i$ of degree less than $m$ can be written in powers of $1-\alpha x$ (substitute
$x=\frac{1-(1-\alpha x)}\alpha$ and expand): $P_i=\sum_{j=0}^{m-1}c_j(1-\alpha x)^j$. So
$\frac{P_i}{(1-\alpha x)^m}=\sum_{j=1}^m\frac{c_{m-j}}{(1-\alpha x)^j}$, and by Proposition~\ref{prop:inverse} its coefficient
of $x^n$ is $\alpha^n\sum_jc_{m-j}\binom{n+j-1}{j-1}$, a polynomial in $n$ of degree less than $m$ times $\alpha^n$.
\end{proof}

\begin{example}[Fibonacci]\label{ex:fibonacci}
$F_0=0$, $F_1=1$, $F_{n+2}=F_{n+1}+F_n$. Here $Q=1-x-x^2$ and $P=QU$ has coefficients $F_0=0$ and $F_1-F_0=1$, so
$U=\frac x{1-x-x^2}$. With $\varphi=\frac{1+\sqrt5}2$ and $\psi=\frac{1-\sqrt5}2$, $1-x-x^2=(1-\varphi x)(1-\psi x)$, and
\[
\begin{gathered}\frac x{(1-\varphi x)(1-\psi x)}=\frac1{\sqrt5}\Bigl(\frac1{1-\varphi x}-\frac1{1-\psi x}\Bigr),\\F_n=\frac{\varphi^n-\psi^n}{\sqrt5}.\end{gathered}
\]
This is Binet's formula; Exercise~\ref{ex:binet} draws consequences.
\end{example}

The method also handles recurrences that are not of constant-coefficient type, as long as the generating
function satisfies an equation one can solve. Two examples from the analysis of algorithms show the range.

\begin{example}[merge sort]\label{ex:mergesort}
Merge sort splits a list of $n=2^m$ items in two halves, sorts each, and merges them with at most $n$ comparisons. Its
worst-case cost $T(n)$ satisfies $T(n)=2T(\frac n2)+n$, so $t_m=T(2^m)$ satisfies $t_m=2t_{m-1}+2^m$. Multiplying by
$y^m$ and summing over $m\ge1$, the series $\tau(y)=\sum t_my^m$ satisfies $\tau-t_0=2y\tau+\frac{2y}{1-2y}$, so
\[
\tau(y)=\frac{t_0}{1-2y}+\frac{2y}{(1-2y)^2},\qq t_m=t_02^m+m2^m,
\]
by Proposition~\ref{prop:inverse}. In terms of $n$: $T(n)=n\log_2n+T(1)n$.
\end{example}

\begin{example}[quicksort]\label{ex:quicksortexact}
Randomised quicksort (Tome~3, Exercise~2.3) chooses a uniform pivot among $n$ items, compares it with the $n-1$
others, and recurses on both sides. Its expected number of comparisons satisfies $C_0=0$ and
\[
C_n=n-1+\frac2n\sum_{k=0}^{n-1}C_k\qquad(n\ge1).
\]
Multiply by $nx^n$ and sum: with $C(x)=\sum C_nx^n$, and $\sum_nn(n-1)x^n=\frac{2x^2}{(1-x)^3}$ (differentiate
$\frac1{1-x}$ twice), this becomes $xC'=\frac{2x^2}{(1-x)^3}+\frac{2xC}{1-x}$, a linear differential equation. Dividing
by $x$, the series $(1-x)^2C$ has derivative $(1-x)^2C'-2(1-x)C=\frac{2x}{1-x}$, so, both sides vanishing at $0$,
\[
\begin{gathered}(1-x)^2C=2\bigl(-\log(1-x)-x\bigr),\\C=\frac2{(1-x)^2}\log\frac1{1-x}-\frac{2x}{(1-x)^2}.\end{gathered}
\]
Since $[x^n]\frac1{(1-x)^2}\log\frac1{1-x}=\sum_{k=1}^n\frac{n+1-k}k=(n+1)H_n-n$ and $[x^n]\frac x{(1-x)^2}=n$,
\[
\begin{aligned}C_n&=2(n+1)H_n-4n\eqbr=2n\ln n+(2\gamma-4)n+O(\ln n),\end{aligned}
\]
using $H_n=\ln n+\gamma+O(\frac1n)$ (Tome~0, Chapter~5, where $\gamma\approx0.5772$ is Euler's constant).
\end{example}

\section{Exponential generating functions}

When the objects counted are built on a set of labels $\{1,\dots,n\}$ (permutations, set partitions, graphs), the right
generating function is the \emph{exponential} one, $\hat A(x)=\sum_na_n\frac{x^n}{n!}$. The reason is the binomial
coefficient in its product:
\[
\hat A\hat B=\sum_n\Bigl(\sum_{k=0}^n\binom nka_kb_{n-k}\Bigr)\frac{x^n}{n!}.
\]
If $a_k$ counts the structures of a kind $\mathcal A$ on a $k$-element set and $b_k$ those of a kind $\mathcal B$, then
$\sum_k\binom nka_kb_{n-k}$ counts the ways to split an $n$-set into two parts $S$, $T$ (with $S$ first) and put an
$\mathcal A$-structure on $S$ and a $\mathcal B$-structure on $T$. Only the size of a set matters for the number of
structures on it, which is what ``labelled'' means.

\begin{theorem}[exponential formula]\label{thm:expformula}
Let $c_k$ be the number of ``connected'' structures of some kind on a $k$-element set ($k\ge1$), and let $h_n$ be the
number of ways to partition an $n$-element set into blocks and put a connected structure on each block ($h_0=1$).
Then, with $\hat C=\sum_{k\ge1}c_k\frac{x^k}{k!}$,
\[
\sum_{n\ge0}h_n\frac{x^n}{n!}=\exp\bigl(\hat C(x)\bigr).
\]
\end{theorem}
\begin{proof}
By the product rule for exponential series, iterated, $[\frac{x^n}{n!}]\hat C^k$ is the number of ways to choose an
\emph{ordered} sequence $(S_1,\dots,S_k)$ of non-empty disjoint sets with union $\{1,\dots,n\}$ and a connected
structure on each. The blocks of a set partition are distinct non-empty sets, so each unordered choice with $k$
blocks appears exactly $k!$ times in this count. Hence the number of structures with exactly $k$ blocks is
$[\frac{x^n}{n!}]\frac{\hat C^k}{k!}$, and summing over $k$ gives $[\frac{x^n}{n!}]\exp(\hat C)$.
\end{proof}

\begin{example}\label{ex:expformula}
\begin{enumerate}[label=(\alph*)]
\item \emph{Permutations.} A permutation is a set of cycles, and there are $(k-1)!$ cyclic orders on a $k$-set. So
$\hat C=\sum_k\frac{x^k}k=\log\frac1{1-x}$ and $\exp(\hat C)=\frac1{1-x}$ by Proposition~\ref{prop:explog}: there are $n!$
permutations, as there should be.
\item \emph{Derangements}, permutations without fixed points: cycles of length at least $2$, so
$\hat D=\exp\bigl(\log\frac1{1-x}-x\bigr)=\frac{e^{-x}}{1-x}$, and $D_n=n!\sum_{k=0}^n\frac{(-1)^k}{k!}$, which is the integer
nearest to $\frac{n!}e$ for $n\ge1$.
\item \emph{Set partitions.} Each block carries one structure, so $c_k=1$ and $\hat C=e^x-1$. The Bell numbers $b_n$,
which count all partitions of an $n$-set, have $\sum b_n\frac{x^n}{n!}=e^{e^x-1}$; they are $1,1,2,5,15,52,203,877,\dots$
Marking each block with a variable $u$ gives $\sum_{n,k}\stirlingii nk u^k\frac{x^n}{n!}=e^{u(e^x-1)}$, the generating
function of the Stirling numbers of Chapter~1.
\item \emph{Involutions}, permutations with $\sigma^2=\mathrm{id}$: cycles of length $1$ or $2$, so their exponential
generating function is $e^{x+x^2/2}$ (Exercise~\ref{ex:involutions}).
\end{enumerate}
\end{example}

The formula also runs backwards: knowing all structures, the logarithm counts the connected ones. Exercise
\ref{ex:connected} counts connected graphs this way, a problem with no simple direct solution.

\section{Partitions of integers}

A \emph{partition} of $n\ge0$ is a way of writing $n$ as a sum of positive integers, the order of the terms not
mattering: $4=3+1=2+2=2+1+1=1+1+1+1$, so $p(4)=5$. We write a partition as $\lambda_1\ge\lambda_2\ge\dots\ge\lambda_r$, and draw
its \emph{Ferrers diagram}: $r$ rows of dots, of lengths $\lambda_1,\dots,\lambda_r$, left aligned (Figure~\ref{fig:ferrers}).
Reflecting the diagram in its diagonal gives the \emph{conjugate} partition, whose parts are the column lengths.

\begin{figure}[ht]
\centering
\fig{ch02-ferrers}
\caption{The Ferrers diagram of $6+4+4+2+1$ and of its conjugate $5+4+3+3+1+1$. Conjugation shows that the number of
partitions of $n$ into at most $k$ parts equals the number of partitions of $n$ into parts at most $k$.}
\label{fig:ferrers}
\end{figure}

\begin{proposition}[Euler]\label{prop:eulerproduct}
In $\Z[[x]]$,
\[
\begin{gathered}\sum_{n\ge0}p(n)x^n=\prod_{k\ge1}\frac1{1-x^k},\qq\prod_{k\ge1}(1+x^k)=\prod_{k\ge1}\frac1{1-x^{2k-1}}.\end{gathered}
\]
The second identity says that the partitions of $n$ into distinct parts are as many as the partitions of $n$ into odd
parts.
\end{proposition}
\begin{proof}
The products make sense since $\frac1{1-x^k}-1$ has order $k$. Expanding $\prod_k(1+x^k+x^{2k}+\cdots)$, the coefficient of
$x^n$ counts the ways to choose $m_k\ge0$ for each $k$, almost all zero, with $\sum_km_kk=n$: a partition with $m_k$ parts
equal to $k$. Similarly $\prod(1+x^k)$ counts partitions into distinct parts, and $\prod\frac1{1-x^{2k-1}}$ partitions into odd
parts. For the identity, $1+x^k=\frac{1-x^{2k}}{1-x^k}$, so for each $N$,
\[
\begin{aligned}\prod_{k\le2N}(1+x^k)&=\frac{\prod_{k\le2N}(1-x^{2k})}{\prod_{k\le2N}(1-x^k)}\eqbr=\frac{\prod_{N<k\le2N}(1-x^{2k})}{\prod_{k\le N}(1-x^{2k-1})}.\end{aligned}
\]
Comparing coefficients of $x^n$ for $n<2N$, where the factors $1-x^{2k}$ with $k>N$ and $1-x^m$ with $m\ge2N$ do not
matter, gives the identity.
\end{proof}

Now the surprise. Expand $\prod_k(1-x^k)$, the inverse of the partition series:
\[
\begin{aligned}\prod_{k\ge1}(1-x^k)&=1-x-x^2+x^5+x^7\eqbr\quad-x^{12}-x^{15}+x^{22}+x^{26}-\cdots\end{aligned}
\]
Almost all coefficients vanish, the others are $\pm1$, and the exponents $1,2,5,7,12,15,22,26,\dots$ are the
\emph{pentagonal numbers} $\frac{j(3j-1)}2$ for $j=\pm1,\pm2,\dots$

\begin{theorem}[Euler's pentagonal number theorem]\label{thm:pentagonal}
\[
\begin{aligned}\prod_{k\ge1}(1-x^k)&=\sum_{j\in\Z}(-1)^jx^{j(3j-1)/2}\\&=1+\sum_{r\ge1}(-1)^rx^{r(3r-1)/2}\eqbr\quad+\sum_{r\ge1}(-1)^rx^{r(3r+1)/2}.\end{aligned}
\]
\end{theorem}
\begin{proof}[Proof (Franklin, 1881)]
Expanding the product, each partition of $n$ into distinct parts contributes $(-1)^r$, where $r$ is its number of
parts. So $[x^n]\prod(1-x^k)=e_n-o_n$, where $e_n$ and $o_n$ count the partitions of $n$ into an even, respectively odd,
number of distinct parts. We pair most of them off.

For a partition $\lambda_1>\dots>\lambda_r$ into distinct parts ($r\ge1$), let $s=\lambda_r$ be its smallest part and let its
\emph{slope} $\sigma$ be the largest $i$ such that $\lambda_1,\dots,\lambda_i$ are consecutive integers, $\lambda_j=\lambda_1-j+1$ for
$j\le i$ (Figure~\ref{fig:franklin}). Call $\lambda$ \emph{exceptional} if $\sigma=r$ and $s\in\{\sigma,\sigma+1\}$. Define:
\begin{itemize}
\item move $A$, when $s\le\sigma$: remove the smallest part and add $1$ to each of the $s$ largest parts;
\item move $B$, when $s>\sigma$: subtract $1$ from each of the $\sigma$ largest parts and add a new part equal to $\sigma$.
\end{itemize}
Let $\mathcal A$ be the non-exceptional partitions with $s\le\sigma$ and $\mathcal B$ those with $s>\sigma$.

\emph{$A$ maps $\mathcal A$ to $\mathcal B$.} Let $\lambda\in\mathcal A$. Then $s\le r-1$: either $\sigma<r$ and $s\le\sigma$, or $\sigma=r$
and $s<\sigma$ since $\lambda$ is not exceptional. So after removing $\lambda_r$ there are at least $s$ parts left, and the result
$\mu$ has parts $\lambda_1+1>\dots>\lambda_s+1>\lambda_{s+1}>\dots>\lambda_{r-1}$, which are distinct since $\lambda_s+1>\lambda_{s+1}$. Its $s$
largest parts are consecutive, and the next one (if any) satisfies $(\lambda_s+1)-\lambda_{s+1}\ge2$: if $s<\sigma$ because
$\lambda_{s+1}=\lambda_s-1$, and if $s=\sigma$ because the slope of $\lambda$ stops at $\sigma$. So $\sigma(\mu)=s$, while the smallest part
of $\mu$ is at least $\lambda_{r-1}>s$. Hence $s(\mu)>\sigma(\mu)$. And $\mu$ is not exceptional: that would need $\sigma(\mu)=r-1$,
so $s=r-1$, and $s(\mu)=\lambda_{r-1}+1=\sigma(\mu)+1=r$, so $\lambda_{r-1}=r-1=s=\lambda_r$, which is impossible.

\emph{$B$ maps $\mathcal B$ to $\mathcal A$.} Let $\mu\in\mathcal B$ with $r$ parts, slope $\sigma$ and smallest part $s>\sigma$. The
result $\nu$ has parts $\mu_1-1>\dots>\mu_\sigma-1$, then $\mu_{\sigma+1},\dots,\mu_r$ if $\sigma<r$, then $\sigma$. These are distinct: if
$\sigma<r$, then $\mu_{\sigma+1}\le\mu_\sigma-2$ by maximality of the slope, and $\mu_r=s>\sigma$; if $\sigma=r$, then $\mu_r-1=s-1>\sigma$
because $s\ne\sigma+1$ ($\mu$ is not exceptional). So $\nu$ is a partition into $r+1$ distinct parts, with smallest part $\sigma$
and slope at least $\sigma$, since its $\sigma$ largest parts are consecutive. So $s(\nu)\le\sigma(\nu)$. If $\nu$ were exceptional
we would have $\sigma(\nu)=r+1$ and $s(\nu)\in\{\sigma(\nu),\sigma(\nu)+1\}$, so $\sigma\ge r+1$, which is impossible. So $\nu\in\mathcal A$.

\emph{$A$ and $B$ are inverse.} For $\lambda\in\mathcal A$, the partition $\mu=A(\lambda)$ has slope $s$, so $B$ subtracts $1$ from its
$s$ largest parts and adds a part $s$: this is $\lambda$. For $\mu\in\mathcal B$, the partition $\nu=B(\mu)$ has smallest part $\sigma$,
so $A$ removes it and adds $1$ to the $\sigma$ largest parts: this is $\mu$.

So $A$ is a bijection from $\mathcal A$ onto $\mathcal B$ that changes the number of parts by one, and the non-exceptional
partitions cancel in $e_n-o_n$. An exceptional partition with $r$ parts is $(2r-1)+(2r-2)+\dots+r$, of sum $\frac{r(3r-1)}2$
(when $s=\sigma=r$), or $2r+(2r-1)+\dots+(r+1)$, of sum $\frac{r(3r+1)}2$ (when $s=r+1$). These numbers are all distinct for
$r\ge1$, so each $n\ge1$ has at most one exceptional partition, contributing $(-1)^r$, and $n=0$ has the empty partition,
contributing $1$. Writing $j=-r$ in the second family, $\frac{r(3r+1)}2=\frac{j(3j-1)}2$, gives the sum over $j\in\Z$.
\end{proof}

\begin{figure}[ht]
\centering
\fig{ch02-franklin}
\caption{Franklin's involution. Left: $7+6+5+2$, with slope $\sigma=3$ (orange) and smallest part $s=2$ (blue). Since
$s\le\sigma$, move $A$ puts the smallest part onto the slope: $8+7+5$, whose slope is $2$ and smallest part $5$. Move $B$
takes the slope back off. The two partitions have $4$ and $3$ parts, so they cancel in $\prod(1-x^k)$.}
\label{fig:franklin}
\end{figure}

\begin{corollary}[Euler's recurrence]\label{cor:partitionrec}
For $n\ge1$, with $p(m)=0$ for $m<0$,
\[
\begin{aligned}p(n)&=p(n-1)+p(n-2)\eqbr\quad-p(n-5)-p(n-7)\eqbr\quad+p(n-12)+p(n-15)-\cdots\end{aligned}
\]
the signs following the pattern $++--$ and the arguments $n-\frac{j(3j-1)}2$ for $j=1,-1,2,-2,\dots$
\end{corollary}
\begin{proof}
Multiply $\sum p(n)x^n$ by $\prod(1-x^k)$: the product is $1$, so for $n\ge1$,
$\sum_{j\in\Z}(-1)^jp\bigl(n-\frac{j(3j-1)}2\bigr)=0$. The term $j=0$ is $p(n)$; move the others to the right.
\end{proof}

This is how MacMahon computed $p(n)$ up to $n=200$ by hand in 1918, finding $p(200)=3\,972\,999\,029\,388$; the
numbers grow like $e^{\pi\sqrt{2n/3}}$, as Chapter~11 will explain, and the recurrence needs only about $\sqrt n$ terms
for each new value.

\section{Exercises}

\begin{exercise}[$\star$]\label{ex:binet}
Deduce from Binet's formula that $F_n$ is the integer nearest to $\frac{\varphi^n}{\sqrt5}$, that
$\frac{F_{n+1}}{F_n}\to\varphi$, and that $F_{n-1}F_{n+1}-F_n^2=(-1)^n$.
\end{exercise}

\begin{exercise}[$\star$]\label{ex:change}
In how many ways can one pay $1$ euro with coins of $1,2,5,10,20,50$ cents and $1$ euro? Explain why multiplying a
series by $\frac1{1-x^c}$ is a running sum with step $c$, and turn this into an algorithm (the Tome~0 notebook asked for
this number).
\end{exercise}

\begin{exercise}[$\star\star$]\label{ex:dobinski}
Show that the Bell numbers satisfy $b_n=\frac1e\sum_{k\ge0}\frac{k^n}{k!}$ (Dobi\'nski's formula), and that
$b_{n+1}=\sum_{k=0}^n\binom nkb_k$.
\end{exercise}

\begin{exercise}[$\star\star$]\label{ex:involutions}
Let $t_n$ be the number of involutions of an $n$-set. Show that $t_{n+1}=t_n+nt_{n-1}$ in two ways: from the generating
function $e^{x+x^2/2}$, and directly.
\end{exercise}

\begin{exercise}[$\star\star$]\label{ex:connected}
There are $2^{\binom n2}$ graphs on the vertex set $\{1,\dots,n\}$. Show that the numbers $c_n$ of \emph{connected}
ones satisfy $\sum_{n\ge1}c_n\frac{x^n}{n!}=\log\sum_{n\ge0}2^{\binom n2}\frac{x^n}{n!}$, and compute $c_1,\dots,c_5$.
\end{exercise}

\begin{exercise}[$\star$]\label{ex:motzkin}
A \emph{Motzkin path} of length $n$ goes from $(0,0)$ to $(n,0)$ with steps $(1,1)$, $(1,-1)$, $(1,0)$, never below the
axis. Show that their generating function satisfies $M=1+xM+x^2M^2$, deduce $M=\frac{1-x-\sqrt{1-2x-3x^2}}{2x^2}$, and
compute the first values $1,1,2,4,9,21,51,\dots$
\end{exercise}

\begin{exercise}[$\star\star\star$]\label{ex:sigma}
Let $\sigma(n)$ be the sum of the divisors of $n$. Taking the logarithmic derivative of the pentagonal theorem, show
that for $n\ge1$
\[
\begin{aligned}\sigma(n)&=\sigma(n-1)+\sigma(n-2)\eqbr\quad-\sigma(n-5)-\sigma(n-7)\eqbr\quad+\sigma(n-12)+\sigma(n-15)-\cdots,\end{aligned}
\]
with the same signs as in Corollary~\ref{cor:partitionrec}, the convention $\sigma(0)=n$ for the term $\sigma(n-n)$
if $n$ is pentagonal, and $\sigma(m)=0$ for $m<0$. Euler found this in 1747 and called it ``a most extraordinary law of the
numbers relating to the sum of their divisors''.
\end{exercise}

\section{Hints}

\begin{hint}{ex:binet}
$|\psi|<1$, so $\bigl|\frac{\psi^n}{\sqrt5}\bigr|<\frac12$. For the last identity, use $\varphi\psi=-1$.
\end{hint}
\begin{hint}{ex:change}
If $B=\frac A{1-x^c}$ then $B=A+x^cB$, that is, $b_n=a_n+b_{n-c}$.
\end{hint}
\begin{hint}{ex:dobinski}
$e^{e^x}=\sum_k\frac{e^{kx}}{k!}$, a double series with non-negative terms for $x\ge0$. For the recurrence, differentiate
$e^{e^x-1}$.
\end{hint}
\begin{hint}{ex:involutions}
Differentiate: $T'=(1+x)T$. Directly: in an involution of $\{1,\dots,n+1\}$, the element $n+1$ is fixed or swapped with one
of the others.
\end{hint}
\begin{hint}{ex:connected}
A graph is a set partition of its vertices into connected components, with a connected graph on each.
\end{hint}
\begin{hint}{ex:motzkin}
Look at the first step, and if it goes up, at the first return to the axis.
\end{hint}
\begin{hint}{ex:sigma}
$x\,D\log\prod(1-x^k)=-\sum_k\frac{kx^k}{1-x^k}=-\sum_{n\ge1}\sigma(n)x^n$. Multiply by $\prod(1-x^k)$ and use the theorem twice.
\end{hint}

\section{Solutions}

\begin{solution}{ex:binet}
\idea The term $\psi^n$ is a small correction that decays geometrically.
\writeup Since $|\psi|=\frac{\sqrt5-1}2<1$, $\bigl|F_n-\frac{\varphi^n}{\sqrt5}\bigr|=\frac{|\psi|^n}{\sqrt5}\le\frac1{\sqrt5}<\frac12$, so $F_n$ is the
integer nearest to $\frac{\varphi^n}{\sqrt5}$. Next,
$\frac{F_{n+1}}{F_n}=\frac{\varphi^{n+1}-\psi^{n+1}}{\varphi^n-\psi^n}=\varphi\cdot\frac{1-(\psi/\varphi)^{n+1}}{1-(\psi/\varphi)^n}\to\varphi$ since
$|\psi/\varphi|<1$. Finally, with $\varphi\psi=-1$ and $\varphi^2+\psi^2=(\varphi+\psi)^2-2\varphi\psi=1+2=3$,
\[
\begin{aligned}5F_{n-1}F_{n+1}&=\varphi^{2n}+\psi^{2n}-(\varphi\psi)^{n-1}(\varphi^2+\psi^2)\eqbr=\varphi^{2n}+\psi^{2n}+3(-1)^n,\\
5F_n^2&=\varphi^{2n}+\psi^{2n}-2(\varphi\psi)^n\eqbr=\varphi^{2n}+\psi^{2n}-2(-1)^n,\end{aligned}
\]
and subtracting, $5(F_{n-1}F_{n+1}-F_n^2)=5(-1)^n$. This is Cassini's identity.
\end{solution}

\begin{solution}{ex:change}
\idea The generating function is a product of one factor $\frac1{1-x^c}$ per coin, and each factor is cheap to apply.
\writeup Let $V=\{1,2,5,10,20,50,100\}$. The number of ways is $[x^{100}]\prod_{c\in V}\frac1{1-x^c}$, by the argument of
Proposition~\ref{prop:eulerproduct}: a choice of $m_c$ coins of value $c$ with $\sum cm_c=100$. If $B=\frac A{1-x^c}$ then
$(1-x^c)B=A$, that is, $b_n=a_n+b_{n-c}$ (with $b_n=0$ for $n<0$): a running sum with step $c$. Starting from
$a=(1,0,0,\dots)$ and applying the seven factors in turn, each in $101$ additions, gives $4563$. This is the classical
dynamic programming algorithm for the coin-change problem, and the generating function explains why the order of the
coins does not matter.
\end{solution}

\begin{solution}{ex:dobinski}
\idea Expand $e^{e^x}$ by the exponential series in the outer variable, then each $e^{kx}$ in $x$.
\writeup For real $x$, $e^{e^x}=\sum_k\frac{e^{kx}}{k!}=\sum_k\sum_n\frac{k^nx^n}{k!\,n!}$. For $x\ge0$ the terms are non-negative,
so the order of summation can be exchanged by Lemma~\ref{lem:packets}; for $x<0$ the double series is summable, being
dominated by the case $|x|$, and Lemma~\ref{lem:packets}(c) applies. Hence $e^{e^x-1}=\sum_n\bigl(\frac1e\sum_k\frac{k^n}{k!}\bigr)\frac{x^n}{n!}$ for every
real $x$, a power series converging on $\R$. By the remark at the end of the first section its coefficients are those of the
formal series $e^{e^x-1}=\sum b_n\frac{x^n}{n!}$, which gives Dobi\'nski's formula. For the recurrence,
$D\,e^{e^x-1}=e^x\,e^{e^x-1}$; the coefficient of $\frac{x^n}{n!}$ on the left is $b_{n+1}$ and on the right
$\sum_k\binom nkb_k\cdot1$, by the product rule for exponential series.
\end{solution}

\begin{solution}{ex:involutions}
\idea A differential equation for the generating function gives a recurrence for the coefficients.
\writeup With $T=e^{x+x^2/2}$, the chain rule gives $T'=(1+x)T$. The coefficient of $\frac{x^n}{n!}$ in $T'$ is $t_{n+1}$; in $T$ it
is $t_n$, and in $xT=\sum t_m\frac{x^{m+1}}{m!}$ it is $n\,t_{n-1}$. So $t_{n+1}=t_n+nt_{n-1}$. Directly: in an involution of
$\{1,\dots,n+1\}$, either $n+1$ is fixed, and the rest is any involution of $\{1,\dots,n\}$, or $n+1$ is swapped with one of
the $n$ others, $i$, and the rest is an involution of the remaining $n-1$ elements. The values are $1,1,2,4,10,26,76,232$.
\end{solution}

\begin{solution}{ex:connected}
\idea All graphs are assemblies of connected ones, so the exponential formula applies and its inverse is the logarithm.
\writeup A graph on $\{1,\dots,n\}$ is the same as the partition of the vertex set into its connected components, with a
connected graph on each block. By Theorem~\ref{thm:expformula}, $\sum_n2^{\binom n2}\frac{x^n}{n!}=\exp(\hat C)$, and by
Proposition~\ref{prop:explog} (third identity, with $1+F$ the series of all graphs), $\hat C=\log\sum_n2^{\binom n2}\frac{x^n}{n!}$.
Expanding $\log(1+F)=F-\frac{F^2}2+\frac{F^3}3-\cdots$ with $F=x+\frac{2x^2}2+\frac{8x^3}6+\frac{64x^4}{24}+\frac{1024x^5}{120}$ to order $5$ gives
\[
\hat C=x+\frac{x^2}2+\frac{2x^3}3+\frac{19x^4}{12}+\frac{91x^5}{15}+\cdots,
\]
so $c_1,\dots,c_5=1,1,4,38,728$. Check $c_3=4$: the path in three ways, and the triangle. Almost all graphs are connected:
$\frac{c_n}{2^{\binom n2}}\to1$, as the notebook shows.
\end{solution}

\begin{solution}{ex:motzkin}
\idea Decompose a path at its first step, and an up step at its matching down step, exactly as for Catalan numbers.
\writeup A non-empty Motzkin path either starts with a flat step followed by any Motzkin path, or starts with an up step;
in the second case, let the first return to the axis be at step $k+2$: between the up step and its matching down step
lies a Motzkin path of length $k$ (raised by one), and after it any Motzkin path. So $M=1+xM+x^2M^2$. Solving the quadratic
$x^2M^2+(x-1)M+1=0$ gives $M=\frac{1-x\pm\sqrt{(1-x)^2-4x^2}}{2x^2}$, and only the minus sign gives a power series (the plus
sign has a pole of order $2$ at $0$). Here $\sqrt{1-2x-3x^2}$ means the formal series $\exp\bigl(\frac12\log(1-2x-3x^2)\bigr)$,
whose square is $1-2x-3x^2$ by Proposition~\ref{prop:explog}. The recurrence $M_n=M_{n-1}+\sum_{k=0}^{n-2}M_kM_{n-2-k}$, read off
$M=1+xM+x^2M^2$, gives $1,1,2,4,9,21,51,127,323$.
\end{solution}

\begin{solution}{ex:sigma}
\idea A logarithmic derivative turns a product into a sum; for $\prod(1-x^k)$ the sum is a divisor sum, and for the
pentagonal series it is a sparse sum.
\writeup Let $E=\prod(1-x^k)$. The logarithm of an infinite product of series $1+A_k$ (with orders of $A_k$ tending to
infinity) is the sum of the logarithms, by Proposition~\ref{prop:explog} applied to finite products and comparison of
coefficients. So $\log E=\sum_k\log(1-x^k)$ and
\[
\begin{aligned}x\frac{DE}E&=\sum_kx\frac{-kx^{k-1}}{1-x^k}\eqbr=-\sum_k\sum_{m\ge1}kx^{km}\eqbr=-\sum_{n\ge1}\sigma(n)x^n,\end{aligned}
\]
grouping the terms with $km=n$. Hence $-E\sum\sigma(n)x^n=xDE$. By Theorem~\ref{thm:pentagonal}, $E=\sum_j(-1)^jx^{g_j}$ and
$xDE=\sum_j(-1)^jg_jx^{g_j}$, with $g_j=\frac{j(3j-1)}2$. Comparing coefficients of $x^n$:
\[
-\sum_j(-1)^j\sigma(n-g_j)=(-1)^{j_0}n\ \text{ if }n=g_{j_0},\ \text{else }0,
\]
with $\sigma(m)=0$ for $m\le0$. The term $j=0$ is $\sigma(n)$, so $\sigma(n)=\sum_{j\ne0}(-1)^{j+1}\sigma(n-g_j)$, plus $-(-1)^{j_0}n$
when $n=g_{j_0}$; this is the convention ``$\sigma(0)=n$'' in the term $-(-1)^{j_0}\sigma(n-g_{j_0})$. For instance $\sigma(5)=\sigma(4)+\sigma(3)-5=7+4-5=6$.
\end{solution}

\section*{Sources and further reading}

\begin{itemize}
\item H.~Wilf, \emph{generatingfunctionology}, 3rd ed., A K Peters, 2006 (free online). The source of the epigraph and of
the exponential formula in this form.
\item R.~Graham, D.~Knuth and O.~Patashnik, \emph{Concrete Mathematics}, 2nd ed., Addison-Wesley, 1994, Chapter~7.
\item P.~Flajolet and R.~Sedgewick, \emph{Analytic Combinatorics}, Cambridge University Press, 2009, Part~A (labelled and
unlabelled constructions, free online).
\item R.~Stanley, \emph{Enumerative Combinatorics}, Vol.~1, 2nd ed., Cambridge University Press, 2012, Chapter~1, and
Vol.~2, 1999, Chapter~5 (the exponential formula).
\item G.~Andrews, \emph{The Theory of Partitions}, Addison-Wesley, 1976, Chapter~1; G.~Andrews and K.~Eriksson,
\emph{Integer Partitions}, Cambridge University Press, 2004.
\item F.~Franklin, \emph{Sur le d\'eveloppement du produit infini $(1-x)(1-x^2)(1-x^3)\cdots$}, Comptes Rendus de
l'Acad\'emie des Sciences 92 (1881), 448--450.
\item G.~H.~Hardy and E.~M.~Wright, \emph{An Introduction to the Theory of Numbers}, 6th ed., Oxford University Press,
2008, Chapter~19.
\item D.~Knuth, \emph{The Art of Computer Programming}, Vol.~3, 2nd ed., Addison-Wesley, 1998, Section~5.2.2 (the analysis of
quicksort).
\end{itemize}
